\documentclass[11pt]{amsart}
\usepackage{amsmath,amsthm,amssymb}
\usepackage{mathtools}
\usepackage{tikz-cd}
\usepackage[backend=biber,style=alphabetic]{biblatex}
\usepackage{hyperref}
\addbibresource{refs.bib}

\newtheorem{theorem}{Theorem}[section]
\newtheorem{lemma}[theorem]{Lemma}
\newtheorem{proposition}[theorem]{Proposition}
\newtheorem{corollary}[theorem]{Corollary}
\theoremstyle{definition}
\newtheorem{definition}[theorem]{Definition}
\newtheorem{example}[theorem]{Example}
\theoremstyle{remark}
\newtheorem{remark}[theorem]{Remark}

\newcommand{\Hom}{\operatorname{Hom}}
\newcommand{\Spec}{\operatorname{Spec}}
\newcommand{\cO}{\mathcal{O}}
\newcommand{\cC}{\mathcal{C}}
\newcommand{\cF}{\mathcal{F}}
\newcommand{\Sh}{\operatorname{Sh}}

\title{Notes on Sheaves and Descent}
\author{A. Mathematician}

\begin{document}
\maketitle

\begin{abstract}
We collect some basic facts on sheaves on sites, descent for quasi-coherent
modules, and the Yoneda lemma. These notes serve as a test document.
\end{abstract}

\tableofcontents

\section{Introduction}\label{sec:intro}
Let $X$ be a scheme and $\cO_X$ its structure sheaf. The theory of descent,
initiated by Grothendieck~\cite{SGA1, is treated at length in~\cite{Vistoli}.
We follow the conventions of~\cite{Stacks}. In Section~\ref{sec:yoneda} we
recall the Yoneda lemma (Lemma~\ref{lem:yoneda}), and in
Section~\ref{sec:descent} we prove Theorem~\ref{thm:fpqc}.

\section{The Yoneda lemma}\label{sec:yoneda}
Let $\cC$ be a locally small category and write $h_X=\Hom_{\cC}(-,X)$.

\begin{lemma}[Yoneda]\label{lem:yoneda}
For every presheaf $F\colon\cC^{\mathrm{op}}\to\mathrm{Set}$ and every object
$X$ of $\cC$ there is a bijection, natural in $X$ and $F$,
\[
\Hom_{\mathrm{PSh}(\cC)}(h_X,F)\xrightarrow{\ \sim\ }F(X),\qquad
\eta\longmapsto\eta_X(\mathrm{id}_X).
\]
\end{lemma}

\begin{proof}
Given $a\in F(X)$ define $\eta^a_Y(g)=F(g)(a)$ for $g\colon Y\to X$.
Naturality of $\eta$ forces the commutative square
\[
\begin{tikzcd}
h_X(X) \arrow[r, "\eta_X"] \arrow[d, "g^*"'] & F(X) \arrow[d, "F(g)"] \\
h_X(Y) \arrow[r, "\eta_Y"'] & F(Y)
\end{tikzcd}
\]
and chasing $\mathrm{id}_X$ gives $\eta_Y(g)=F(g)(\eta_X(\mathrm{id}_X))$.
\end{proof}

\begin{corollary}\label{cor:ff}
The Yoneda embedding $h\colon\cC\to\mathrm{PSh}(\cC)$ is fully faithful.
\end{corollary}

This is used throughout; see~\cite{MacLane}.


\section{Descent}\label{sec:descent}
\begin{definition}\label{def:descent-datum}
Let $f\colon S'\to S$ be a morphism of schemes. A \emph{descent datum} for a
quasi-coherent $\cO_{S'}$-module $\cF'$ is an isomorphism
$\varphi\colon p_1^*\cF'\xrightarrow{\sim}p_2^*\cF'$ on $S''=S'\times_S S'$
satisfying the cocycle condition
\begin{equation}\label{eq:cocycle}
p_{13}^*\varphi = p_{23}^*\varphi\circ p_{12}^*\varphi
\quad\text{on } S'\times_S S'\times_S S'.
\end{equation}
\end{definition}

\begin{theorem}[fpqc descent]\label{thm:fpqc}
Let $f\colon S'\to S$ be faithfully flat and quasi-compact. Then the functor
\[
\mathrm{QCoh}(S)\longrightarrow \mathrm{QCoh}(S'\to S),\qquad
\cF\longmapsto (f^*\cF,\mathrm{can})
\]
is an equivalence of categories.
\end{theorem}

\begin{proof}
By~\cite[Tag 023N]{Stacks} it suffices to treat the affine case
$S=\Spec A$, $S'=\Spec B$ with $A\to B$ faithfully flat. Consider the
Amitsur complex
\begin{equation}\label{eq:amitsur}
0\to M\to M\otimes_A B\rightrightarrows M\otimes_A B\otimes_A B
\to\cdots
\end{equation}
which is exact since it becomes split exact after the faithfully flat base
change $-\otimes_A B$. The cocycle condition~\eqref{eq:cocycle} is exactly
what is needed to descend $\cF'$.
\end{proof}

\begin{remark}
The diagram
\[
\begin{tikzcd}
S'\times_S S' \arrow[r, shift left, "p_1"] \arrow[r, shift right, "p_2"'] & S' \arrow[r, "f"] & S
\end{tikzcd}
\]
is a coequalizer in the category of schemes when $f$ is fpqc.
\end{remark}

\section{Sheaves on sites}\label{sec:sites}
Let $\cC$ be a site. A presheaf $\cF$ on $\cC$ is a sheaf if for every
covering $\{U_i\to U\}_{i\in I}$ the sequence
\[
\cF(U)\to\prod_{i\in I}\cF(U_i)\rightrightarrows\prod_{i,j\in I}\cF(U_i\times_U U_j)
\]
is an equalizer. The category $\Sh(\cC)$ is a Grothendieck topos.

\begin{proposition}\label{prop:sheafification}
The inclusion $\Sh(\cC)\hookrightarrow\mathrm{PSh}(\cC)$ admits an exact left
adjoint $\cF\mapsto\cF^{\#}$.
\end{proposition}

\begin{proof}
Apply the plus construction twice: $\cF^{\#}=(\cF^+)^+$, where
\[
\cF^+(U)=\operatorname*{colim}_{\mathcal{U}\in J(U)}\check{H}^0(\mathcal{U},\cF).
\]
Exactness follows since filtered colimits commute with finite limits.
\end{proof}

\begin{example}
For $X$ a topological space, the small site $\mathrm{Open}(X)$ gives the
usual notion. For a scheme $X$ the étale site $X_{\text{ét}}$ gives
étale sheaves, and
\[
H^i_{\text{ét}}(X,\mu_n)\cong H^i(X(\mathbb{C}),\mathbb{Z}/n)
\]
for $X$ smooth over $\mathbb{C}$ by Artin's comparison theorem.
\end{example}

\subsection{Long computation}
\begin{align}
\Hom_{\Sh(\cC)}(\cF^{\#},\mathcal{G})
&\cong \Hom_{\mathrm{PSh}(\cC)}(\cF,\mathcal{G}) \label{eq:adj1}\\
&\cong \lim_{U\in\cC}\Hom_{\mathrm{Set}}(\cF(U),\mathcal{G}(U)) \label{eq:adj2}\\
&\cong \int_{U\in\cC}\mathcal{G}(U)^{\cF(U)}. \label{eq:adj3}
\end{align}
Equations~\eqref{eq:adj1}--\eqref{eq:adj3} combine to give the end formula.

\subsection{Filler subsection 1}
Lorem ipsum dolor sit amet, consectetur adipiscing elit, sed do eiusmod
tempor incididunt ut labore et dolore magna aliqua. For every $n\ge 1$ we have
\[
\sum_{k=0}^{n}\binom{n}{k}x^k y^{n-k}=(x+y)^n,\qquad
\int_0^1 t^{1}\,dt=\frac{1}{2}.
\]
Ut enim ad minim veniam, quis nostrud exercitation ullamco laboris nisi ut
aliquip ex ea commodo consequat. Duis aute irure dolor in reprehenderit in
voluptate velit esse cillum dolore eu fugiat nulla pariatur. Excepteur sint
occaecat cupidatat non proident, sunt in culpa qui officia deserunt mollit
anim id est laborum. See Theorem~\ref{thm:fpqc} and~\cite{Vistoli}.
Lorem ipsum dolor sit amet, consectetur adipiscing elit, sed do eiusmod
tempor incididunt ut labore et dolore magna aliqua. Ut enim ad minim veniam,
quis nostrud exercitation ullamco laboris nisi ut aliquip ex ea commodo
consequat. Duis aute irure dolor in reprehenderit in voluptate velit esse.

\subsection{Filler subsection 2}
Lorem ipsum dolor sit amet, consectetur adipiscing elit, sed do eiusmod
tempor incididunt ut labore et dolore magna aliqua. For every $n\ge 2$ we have
\[
\sum_{k=0}^{n}\binom{n}{k}x^k y^{n-k}=(x+y)^n,\qquad
\int_0^1 t^{2}\,dt=\frac{1}{3}.
\]
Ut enim ad minim veniam, quis nostrud exercitation ullamco laboris nisi ut
aliquip ex ea commodo consequat. Duis aute irure dolor in reprehenderit in
voluptate velit esse cillum dolore eu fugiat nulla pariatur. Excepteur sint
occaecat cupidatat non proident, sunt in culpa qui officia deserunt mollit
anim id est laborum. See Theorem~\ref{thm:fpqc} and~\cite{Vistoli}.
Lorem ipsum dolor sit amet, consectetur adipiscing elit, sed do eiusmod
tempor incididunt ut labore et dolore magna aliqua. Ut enim ad minim veniam,
quis nostrud exercitation ullamco laboris nisi ut aliquip ex ea commodo
consequat. Duis aute irure dolor in reprehenderit in voluptate velit esse.

\subsection{Filler subsection 3}
Lorem ipsum dolor sit amet, consectetur adipiscing elit, sed do eiusmod
tempor incididunt ut labore et dolore magna aliqua. For every $n\ge 3$ we have
\[
\sum_{k=0}^{n}\binom{n}{k}x^k y^{n-k}=(x+y)^n,\qquad
\int_0^1 t^{3}\,dt=\frac{1}{4}.
\]
Ut enim ad minim veniam, quis nostrud exercitation ullamco laboris nisi ut
aliquip ex ea commodo consequat. Duis aute irure dolor in reprehenderit in
voluptate velit esse cillum dolore eu fugiat nulla pariatur. Excepteur sint
occaecat cupidatat non proident, sunt in culpa qui officia deserunt mollit
anim id est laborum. See Theorem~\ref{thm:fpqc} and~\cite{Vistoli}.
Lorem ipsum dolor sit amet, consectetur adipiscing elit, sed do eiusmod
tempor incididunt ut labore et dolore magna aliqua. Ut enim ad minim veniam,
quis nostrud exercitation ullamco laboris nisi ut aliquip ex ea commodo
consequat. Duis aute irure dolor in reprehenderit in voluptate velit esse.

\subsection{Filler subsection 4}
Lorem ipsum dolor sit amet, consectetur adipiscing elit, sed do eiusmod
tempor incididunt ut labore et dolore magna aliqua. For every $n\ge 4$ we have
\[
\sum_{k=0}^{n}\binom{n}{k}x^k y^{n-k}=(x+y)^n,\qquad
\int_0^1 t^{4}\,dt=\frac{1}{5}.
\]
Ut enim ad minim veniam, quis nostrud exercitation ullamco laboris nisi ut
aliquip ex ea commodo consequat. Duis aute irure dolor in reprehenderit in
voluptate velit esse cillum dolore eu fugiat nulla pariatur. Excepteur sint
occaecat cupidatat non proident, sunt in culpa qui officia deserunt mollit
anim id est laborum. See Theorem~\ref{thm:fpqc} and~\cite{Vistoli}.
Lorem ipsum dolor sit amet, consectetur adipiscing elit, sed do eiusmod
tempor incididunt ut labore et dolore magna aliqua. Ut enim ad minim veniam,
quis nostrud exercitation ullamco laboris nisi ut aliquip ex ea commodo
consequat. Duis aute irure dolor in reprehenderit in voluptate velit esse.

\subsection{Filler subsection 5}
Lorem ipsum dolor sit amet, consectetur adipiscing elit, sed do eiusmod
tempor incididunt ut labore et dolore magna aliqua. For every $n\ge 5$ we have
\[
\sum_{k=0}^{n}\binom{n}{k}x^k y^{n-k}=(x+y)^n,\qquad
\int_0^1 t^{5}\,dt=\frac{1}{6}.
\]
Ut enim ad minim veniam, quis nostrud exercitation ullamco laboris nisi ut
aliquip ex ea commodo consequat. Duis aute irure dolor in reprehenderit in
voluptate velit esse cillum dolore eu fugiat nulla pariatur. Excepteur sint
occaecat cupidatat non proident, sunt in culpa qui officia deserunt mollit
anim id est laborum. See Theorem~\ref{thm:fpqc} and~\cite{Vistoli}.
Lorem ipsum dolor sit amet, consectetur adipiscing elit, sed do eiusmod
tempor incididunt ut labore et dolore magna aliqua. Ut enim ad minim veniam,
quis nostrud exercitation ullamco laboris nisi ut aliquip ex ea commodo
consequat. Duis aute irure dolor in reprehenderit in voluptate velit esse.

\subsection{Filler subsection 6}
Lorem ipsum dolor sit amet, consectetur adipiscing elit, sed do eiusmod
tempor incididunt ut labore et dolore magna aliqua. For every $n\ge 6$ we have
\[
\sum_{k=0}^{n}\binom{n}{k}x^k y^{n-k}=(x+y)^n,\qquad
\int_0^1 t^{6}\,dt=\frac{1}{7}.
\]
Ut enim ad minim veniam, quis nostrud exercitation ullamco laboris nisi ut
aliquip ex ea commodo consequat. Duis aute irure dolor in reprehenderit in
voluptate velit esse cillum dolore eu fugiat nulla pariatur. Excepteur sint
occaecat cupidatat non proident, sunt in culpa qui officia deserunt mollit
anim id est laborum. See Theorem~\ref{thm:fpqc} and~\cite{Vistoli}.
Lorem ipsum dolor sit amet, consectetur adipiscing elit, sed do eiusmod
tempor incididunt ut labore et dolore magna aliqua. Ut enim ad minim veniam,
quis nostrud exercitation ullamco laboris nisi ut aliquip ex ea commodo
consequat. Duis aute irure dolor in reprehenderit in voluptate velit esse.

\printbibliography
\end{document}
